\section{Conclusion}
\label{sec:conclusion}

This paper presented an ambient-aware robust Pyomo-MILP framework for day-ahead scheduling of a heterogeneous multi-chiller plant with thermal storage and auxiliary equipment. The formulation combines unit commitment, piecewise-linear chiller performance, storage dynamics, budgeted robustness for cooling load and electricity price, and a compact wet-bulb-aware treatment of efficiency degradation and available-capacity derating.

On the semi-synthetic 24-hour benchmark, the deterministic MILP with storage achieved the lowest nominal cost but was highly vulnerable to out-of-sample disturbances. The robust formulations sharply reduced expected shortage, and storage partially recovered the associated cost premium. The ambient-aware extension then provided an additional layer of protection under weather-linked stress with only a small increase in nominal cost. These findings support the use of ambient-aware planning-oriented scheduling as a benchmark methodology for full-plant cooling optimization.

The current study has four main limitations. First, the experiments are based on semi-synthetic demand, tariff, and weather profiles rather than measured plant data. Second, the auxiliary-equipment and ambient derating models are intentionally compact and therefore omit richer hydraulic and thermodynamic effects. Third, the work addresses day-ahead scheduling only and does not yet include rolling-horizon feedback or zone-level thermal dynamics. Fourth, the uncertainty model is budgeted and interval-based rather than scenario-driven or distributionally robust.

Future work should proceed along three complementary directions: field-calibrated validation on a measured chiller plant, richer ambient and hydraulic surrogate models that remain optimization-friendly, and tighter integration between robust day-ahead scheduling and real-time supervisory control. Even in its current form, however, the proposed framework provides a transparent and extensible baseline for ambient-aware robust cooling-plant optimization studies.
